\subsection{Electrical subsystem}
\begin{figure}[ht]
    \centering
    \begin{circuitikz}
        \draw
        (0,0) to[V, v=$v_c(t)$] (0,4)
        to[R, l=$R_c$] (2.5,4)
        to[L, l=$L$] (5,4)
        to[R, l=$R_s$] (7.5,4)
        to[short, i=$i(t)$] (7.5,0)
        -- (0,0);
    \end{circuitikz}
    \caption{Electrical model of the electromagnet circuit with coil resistance $R_c$, sense resistance $R_s$, and inductance $L$.}
    \label{fig:rl_detailed}
\end{figure}

The electromagnet is modelled as a resistor-inductor circuit. If $v_c(t)$ is the voltage applied to the coil, the current dynamics are
\begin{equation}
    v_c(t) = R i(t) + L \frac{\dd i(t)}{\dd t}.
    \label{eq:electrical_rl}
\end{equation}

When the amplifier gain is included, the coil voltage can be written as
\begin{equation}
    v_c(t) = K_g u(t),
\end{equation}
where $u(t)$ is the command voltage generated by the controller. Therefore,
\begin{equation}
    \dot{i}(t) = -\frac{R}{L}i(t) + \frac{K_g}{L}u(t).
    \label{eq:current_dynamics}
\end{equation}

In the implemented architecture, an inner current controller is designed first. When this loop is closed, the electrical dynamics can also be approximated by a first-order transfer function
\begin{equation}
    H_i(s) = \frac{I(s)}{I_{ref}(s)} = \frac{\omega_i}{s + \omega_i},
    \label{eq:current_closed_loop}
\end{equation}
where $\omega_i$ is the bandwidth of the closed current loop.
